% POSPSEUDO3
\documentclass[11pt]{article}
\usepackage[margin=.88in]{geometry}
\usepackage{amsmath,amssymb,amsthm}
\newtheorem{theorem}{Theorem}
\begin{document}

\begin{center}
{\Large POSPSEUDO3: Order-three pseudo-independence of moment-balanced UNSAT}\\[4pt]
\textbf{Date: 2026-10-04}\\
\textbf{Status: PROVED/EXACT convex duality + COMPUTATIONAL pilots}\\
\textbf{P versus NP status: OPEN}
\end{center}

\section{Fractional-design formulation}

Let \(\Omega_{q,t}^{UNSAT}\) be the exact-\(t\) UNSAT layer and let
\(U_{q,t}\) denote the uniform exact-\(t\) slice.

For fixed \(k\), let \(\Phi_k(F)\) be the vector of all squarefree clause
monomials through degree \(k\), or equivalently after signed-permutation
averaging, the vector of all canonical orbit counts through order \(k\).

\begin{theorem}[FRACTIONAL-DESIGN-DUALITY]
There exists a probability distribution \(\mu\) supported on
\(\Omega_{q,t}^{UNSAT}\) satisfying
\[
\mathbb E_{F\sim\mu}\Phi_k(F)
=
\mathbb E_{F\sim U_{q,t}}\Phi_k(F)
\]
if and only if the uniform moment vector lies in
\[
\operatorname{conv}
\{\Phi_k(F):F\in\Omega_{q,t}^{UNSAT}\}.
\]

If not, there is a real polynomial \(P\) of degree at most \(k\) whose
expectation separates every UNSAT point from the uniform moment vector.
The separator may be chosen invariant under signed variable permutations.
\end{theorem}

\begin{proof}
The first statement is the definition of membership in the convex hull of the
feature vectors.  If membership fails, finite-dimensional hyperplane
separation gives a linear functional \(a\) and threshold \(b\) separating the
uniform target from every UNSAT feature vector.  Then
\[
P(F)=a\cdot\Phi_k(F)
\]
is a degree-\(k\) polynomial separator in expectation.  Averaging \(P\) over
the signed-variable permutation group preserves its degree and the separation
because both the UNSAT family and the uniform target are invariant.
\end{proof}

This theorem does not by itself turn a separating polynomial into a
constant-gap monotone-circuit distinguisher.  It supplies the exact
moment-feasibility dichotomy.

\section{Order-three local pseudo-independence}

PSEUDOSLICE1 established that for the uniform exact slice, the minimum
\(\varepsilon\) in the TR26-220 pseudo-independence definition can be computed
by a small exact LP.  At \(q=20,t=70\),
\[
\varepsilon^{unif}_3=0.21077.
\]
At \(q=24,t=84\),
\[
\varepsilon^{unif}_3=0.19231.
\]

For an ideal signed-permutation-invariant positive distribution whose
one-coordinate and all pair-coordinate moments exactly match the uniform
slice, every three-coordinate marginal is determined by one remaining number:
the triple inclusion probability \(\theta\).

Thus every canonical triple orbit can be tested by an 8-corner
product-mixture LP.

\section{POS-PSEUDO3 experiment}

The positive distribution was an empirical maximum-entropy tilt of sampled
UNSAT formulas so that all eight independent signed pair-orbit means match the
exact uniform-slice target.

For each retained formula, 2000 uniformly sampled clause triples were
canonicalized.  The control used independent uniform exact-slice formulas with
the same number of sampled triples.

All stages were separately below 30 seconds.

\subsection*{q=20}

Two independent positive pools:

\[
\begin{array}{c|c|c|c}
\text{seed}&\#UNSAT&ESS&\text{pair fit error}\\
\hline
142020&355&53.80&4.6\times10^{-10}\\
152020&321&48.21&4.0\times10^{-11}
\end{array}
\]

For the most common canonical triple orbit, having uniform orbit frequency
approximately \(0.215\), the inferred minimum local pseudo-independence
parameter was
\[
0.270
\qquad\text{and}\qquad
0.272
\]
in the two independent runs.

Among the seven orbit types with uniform frequency at least \(0.03\), the
largest inferred \(\varepsilon_3\) was
\[
0.289
\qquad\text{and}\qquad
0.311.
\]

Thus the common three-coordinate marginals of the positive distribution are
somewhat farther from local product mixtures than the uniform slice itself,
but the excess is modest rather than catastrophic.

\subsection*{q=24}

Two independent positive pools:

\[
\begin{array}{c|c|c}
\text{seed}&\#UNSAT&ESS\\
\hline
162424&258&29.82\\
172424&284&40.36
\end{array}
\]

Here
\[
\varepsilon^{unif}_3=0.1923.
\]

For the six orbit types with uniform frequency at least \(0.03\), the largest
inferred values were
\[
0.313
\qquad\text{and}\qquad
0.298.
\]

The dominant orbit itself gave \(0.200\) and \(0.229\).

There is therefore no clear finite-size evidence yet that the worst common
order-three parameter is shrinking from \(q=20\) to \(q=24\), although the
dominant orbit moves close to the uniform-slice value.

The positive effective sample sizes are only 30--54 and the orbit census is
sampled, so this is not an asymptotic statement.

\section{Sampled full order-three convex fit}

Using the first \(q=20\) pool, 1000 fresh clause triples per UNSAT formula were
used to form an 81-coordinate canonical triple-orbit vector.

A convex combination of the 355 sampled UNSAT vectors could match an
independently sampled uniform target to training max error
\[
1.4\times10^{-5},
\]
using 62 support points and empirical ESS \(38.5\).

However, when the same weights were evaluated with independent resampling of
both the UNSAT formulas and the uniform target, the max error rose to
\[
0.00392.
\]
In fact the original pair-only maximum-entropy weights had a smaller
independent max error of
\[
0.00237.
\]

Therefore the 81-dimensional sampled convex fit is \textbf{overfitting orbit
sampling noise}; it is not evidence that exact order-three fractional design
feasibility has been established.

A lower-dimensional fit using the exact pair coordinates plus the seven most
common triple orbits retained ESS \(51.5\), but its independent triple errors
were again at the \(10^{-3}\) level, comparable with the sampling noise.

\section{Interpretation}

The order-three result is mixed but useful.

\begin{itemize}
\item Pair balancing does not immediately fail at order 3.
\item Common triple marginals can be represented by nearby product mixtures
with \(\varepsilon\) around \(0.2\)--\(0.3\) at these sizes.
\item The worst common value does not yet show a reliable downward trend.
\item High-dimensional sampled moment matching overfits badly enough that an
exact census or stronger regularization is required.
\end{itemize}

This turns the next structural question into a precise one:

\[
\boxed{
\text{Does there exist, for growing }k,
\text{ an UNSAT-supported fractional }k\text{-design on the exact slice?}
}
\]

Equivalently, can the uniform slice moment vector through order \(k\) be kept
inside the UNSAT moment polytope while \(k\) grows?

\section{Next target}

\textbf{FRACTIONAL-DESIGN3-EXACT.}
At \(q=16\) and then \(q=18\), compute exact canonical triple-orbit vectors
for a sufficiently large UNSAT pool in multiple sub-30-second census chunks,
then solve the exact symmetry-reduced convex feasibility problem.  Do not use
sampled triple frequencies in the final feasibility decision.

In parallel, derive quantitative margin consequences of failure of the convex
feasibility problem, and test whether the resulting low-degree separator is a
polynomial-size exact-slice hostile circuit after coefficient-bit accounting.

The 2026 pseudo-independent sunflower theorem remains a candidate downstream
tool only if useful local order can be made to grow and its
\(\varepsilon\le\gamma/w\) condition is met.

\section{Status}

FRACTIONAL-DESIGN-DUALITY: \textbf{PROVED/EXACT}.

POS-PSEUDO3: \textbf{COMPUTATIONAL}.

Exact order-three fractional-design feasibility: \textbf{OPEN}.

Growing-order pseudo-independence of UNSAT-supported distributions:
\textbf{OPEN}.

TWO-DIST-MOMENT-APPROX and the general monotone slice lower bound:
\textbf{OPEN}.

Therefore \(P\ne NP\) remains OPEN.

\end{document}
